% =============================================================================
% Project : THE PHILOSOPHY OF VIABILITY
% Book    : Book I — Ontology of Viability
% Book ID : VS.B01
% Part    : Part I — Identity and Manifestation
% Part ID : VS.B01.P01
% Chapter : Chapter 08 — Comparative Philosophical Examination
% Chapter ID : VS.B01.P01.C08
%
% File :
%   2_MAINMATTER/PART_I_IDENTITY_AND_MANIFESTATION/
%   CHAPTER_08_COMPARATIVE_PHILOSOPHICAL_EXAMINATION/
%   sections/section_10_kant.tex
%
% Section :
%   Section 10 — Immanuel Kant
%
% Structural Identifier :
%   VS.B01.P01.C08.S10
%
% Purpose :
%   To reconstruct Kant's account of the conditions of possible
%   experience, the categories, the transcendental unity of
%   apperception, the paralogisms of rational psychology, the
%   permanence of substance, personal identity, and the
%   distinction between appearances and things in themselves;
%   and to examine the ontological reach and limits of these
%   arguments in relation to the foundations of Part I.
%
% Philosophical Scope :
%   The section concentrates on the ontological implications
%   of Kant's critical philosophy. It does not reconstruct
%   Kant's practical philosophy or treat the categorical
%   imperative as a premise in the present examination.
%
% Foundational Restriction :
%   The conditions under which an object can be experienced
%   or cognized must not be identified without argument with
%   the conditions under which that object exists.
%
%   The unity of apperception is not treated as an ontological
%   creator of the identity of beings. Nor is the identity of
%   apperception equated with the persistence of a substantial
%   soul or with empirical personal identity.
%
% Claim Form and Epistemic Status :
%   Historical reconstructions are textually supported.
%   Interpretive claims are qualified where Kant scholarship
%   remains divided.
%   Comparative arguments are presented as argued assessments.
%   The ontological adequacy of Kant's account remains subject
%   to the distinctions and open questions stated below.
%
% =============================================================================


\section{Immanuel Kant}
\label{sec:ch08_s10}

\index[name]{Kant, Immanuel}
\index[concept]{Transcendental Idealism}
\index[concept]{Copernican Revolution}
\index[concept]{Synthetic A Priori}
\index[concept]{Sensibility}
\index[concept]{Understanding}
\index[concept]{Categories}
\index[concept]{Transcendental Apperception}
\index[concept]{Synthetic Unity of Apperception}
\index[concept]{Self-Consciousness}
\index[concept]{Synthesis}
\index[concept]{Objectivity}
\index[concept]{Rational Psychology}
\index[concept]{Paralogisms}
\index[concept]{Substance}
\index[concept]{Permanence}
\index[concept]{Personal Identity}
\index[concept]{Phenomena}
\index[concept]{Noumena}
\index[concept]{Thing in Itself}
\index[concept]{Transcendental Reflection}


% =============================================================================
% 1. SCOPE AND COMPARATIVE ORIENTATION
% =============================================================================

\subsection*{Scope and Comparative Orientation}

% Structural ID: VS.B01.P01.C08.S10.SS01.P001
\phantomsection\label{para:VS.B01.P01.C08.S10.SS01.P001}

Immanuel Kant is an indispensable interlocutor for an
ontological examination of identity because his critical
philosophy distinguishes the unity required for objective
experience from the empirical succession of representations
and from the metaphysical claims traditionally made about
the soul. His analysis asks what must be presupposed for
a manifold of intuition to be experienced as an ordered
world of objects. It also asks whether the conditions
that make such experience possible authorize conclusions
about the ultimate constitution of the objects experienced.
These questions bear directly on the distinction between
the identity of a being and the conditions under which
that being can be represented as an object.

% Structural ID: VS.B01.P01.C08.S10.SS01.P002
\phantomsection\label{para:VS.B01.P01.C08.S10.SS01.P002}

The comparison requires several distinctions. Kant's
transcendental unity of apperception concerns the unity
under which representations must be combinable in one
self-consciousness. The empirical consciousness of
particular mental states concerns what is given in inner
sense. The persistence of substance concerns the temporal
determination of appearances. Personal identity concerns
the identity of a person through changes of states.
The identity of an object across different presentations
raises yet another problem. These questions are connected
in Kant's critical project, but they are not interchangeable.

% Structural ID: VS.B01.P01.C08.S10.SS01.P003
\phantomsection\label{para:VS.B01.P01.C08.S10.SS01.P003}

The present section examines Kant as a philosopher of the
conditions and limits of objective experience, not as an
author whose theory of cognition can simply be adopted
as an ontology. The central comparative question is whether
his account establishes the conditions under which identity
can be experienced and judged, or whether it also provides
an account of the ontological status of identity itself.
The distinction is essential because an account of the
possibility of identifying an object does not automatically
settle what makes that object the being that it is.

% Structural ID: VS.B01.P01.C08.S10.SS01.P004
\phantomsection\label{para:VS.B01.P01.C08.S10.SS01.P004}

The examination gives priority to the Critique of Pure
Reason. Kant's practical philosophy, including the
categorical imperative, is not developed here because
it does not perform a necessary role in the present
ontological comparison. The discussion of selfhood is
restricted to the distinction between apperceptive unity,
empirical self-consciousness, rational psychology, and
personal identity. This restriction prevents the section
from becoming a general survey of Kant's philosophy.


% =============================================================================
% 2. HISTORICAL AND TEXTUAL LOCATION
% =============================================================================

\subsection{Historical and Textual Location}
\label{subsec:ch08-kant-historical}

% Structural ID: VS.B01.P01.C08.S10.SS02.P001
\phantomsection\label{para:VS.B01.P01.C08.S10.SS02.P001}

Kant's critical philosophy is developed most systematically
in the Critique of Pure Reason, first published in 1781
and substantially revised in 1787. The first edition is
commonly cited by the pagination of the A edition and the
second by that of the B edition. The combined A/B
pagination provides a stable means of locating passages
across the two editions. The present examination uses
these standard references rather than relying on page
numbers that vary among translations.
\cite{KantCritiquePureReason}

% Structural ID: VS.B01.P01.C08.S10.SS02.P002
\phantomsection\label{para:VS.B01.P01.C08.S10.SS02.P002}

The first edition contains the Transcendental Deduction
in a form that differs significantly from the revised
B-edition presentation. The second edition also revises
the Paralogisms and makes several arguments more explicit.
These differences matter for the interpretation of
Kant's account of self-consciousness and the limits of
rational psychology. A reconstruction that cites only
one edition without identifying the relevant textual
context risks presenting a contested interpretation
as though it were the sole formulation of Kant's position.

\footnote{For a detailed examination of the structure
and interpretation of Kant's transcendental idealism,
see Henry E. Allison, \textit{Kant's Transcendental
Idealism: An Interpretation and Defense}, 2nd ed.
(New Haven: Yale University Press, 2004).}

% Structural ID: VS.B01.P01.C08.S10.SS02.P003
\phantomsection\label{para:VS.B01.P01.C08.S10.SS02.P003}

The present discussion concentrates on the Transcendental
Aesthetic, the Transcendental Analytic, the Transcendental
Deduction, the Analogies of Experience, the Paralogisms,
and the Amphiboly of the Concepts of Reflection. These
parts of the Critique address the conditions of intuition,
the unity of experience, the temporal determination of
appearances, the limits of rational psychology, and the
conditions under which identity and difference can be
judged objectively.

% Structural ID: VS.B01.P01.C08.S10.SS02.P004
\phantomsection\label{para:VS.B01.P01.C08.S10.SS02.P004}

Kant's critical project must be understood against the
rationalist and empiricist positions that preceded it.
The relation to Leibniz concerns, among other issues,
the intellectual determination of objects and the
principle of identity of indiscernibles. The relation
to Locke concerns personal identity and the role of
consciousness. Hume's challenge concerning causation,
substance, and the limits of inference provides another
important background. The present comparison therefore
continues the problems examined in the preceding sections
on Leibniz, Locke, and Hume (see \SectionRef{sec:ch08_s09}), while preserving the
differences among their respective arguments.
\cite{KantCritiquePureReason}


% =============================================================================
% 3. THE COPERNICAN TURN AND THE CONDITIONS OF EXPERIENCE
% =============================================================================

\subsection{The Copernican Turn and the Conditions of Experience}
\label{subsec:ch08-kant-copernican-turn}

% Structural ID: VS.B01.P01.C08.S10.SS03.P001
\phantomsection\label{para:VS.B01.P01.C08.S10.SS03.P001}

Kant's Copernican turn reverses the direction of a
traditional explanatory expectation. Instead of assuming
that cognition must conform in every respect to objects
as they are independently given, Kant asks whether
objects, insofar as they can be experienced by us,
must conform to the a priori conditions of our mode
of cognition. The claim is not that the mind creates
the existence of objects through an arbitrary act.
It is that the possibility of objects as objects of
human experience presupposes forms and rules that
belong to the cognitive faculties.
\cite{KantCritiquePureReason}

% Structural ID: VS.B01.P01.C08.S10.SS03.P002
\phantomsection\label{para:VS.B01.P01.C08.S10.SS03.P002}

The critical question is therefore transcendental:
which conditions must obtain for experience to be
possible at all? Kant's answer is not an empirical
generalization about how human beings happen to think.
It concerns conditions that are claimed to be necessary
for objective experience. The argument seeks to explain
how synthetic a priori judgments can possess objective
validity without deriving their necessity from repeated
observation.

% Structural ID: VS.B01.P01.C08.S10.SS03.P003
\phantomsection\label{para:VS.B01.P01.C08.S10.SS03.P003}

This shift changes the status of metaphysical claims.
A concept may be logically coherent while lacking
legitimate application to objects. The fact that reason
can formulate a concept does not establish that an
object corresponding to it is given or knowable.
The critical inquiry therefore distinguishes the
logical use of concepts from their objective use
under conditions that make experience possible.

% Structural ID: VS.B01.P01.C08.S10.SS03.P004
\phantomsection\label{para:VS.B01.P01.C08.S10.SS03.P004}

For the Ontology of Viability, the significance of
this turn lies in the distinction it forces between
the conditions of access to an object and the conditions
of the object's being. Kant's argument establishes
a demanding account of the former. Whether it also
establishes the latter is a further question. The
critical philosophy may constrain what can legitimately
be claimed about identity without thereby demonstrating
that identity is constituted by the conditions of
human experience.

% Structural ID: VS.B01.P01.C08.S10.SS03.P005
\phantomsection\label{para:VS.B01.P01.C08.SS03.P005}

The distinction must not be used to dismiss Kant's
position as merely epistemological. His transcendental
arguments concern the objective validity of categories
and the structure of possible experience. They are
not simply descriptions of subjective habits. The
comparative task is to identify precisely what their
necessity establishes, rather than to reduce the
critical project to an account of individual cognition.


% =============================================================================
% 4. SENSIBILITY, UNDERSTANDING, AND THE CATEGORIES
% =============================================================================

\subsection{Sensibility, Understanding, and the Categories}
\label{subsec:ch08-kant-faculties}

% Structural ID: VS.B01.P01.C08.S10.SS04.P001
\phantomsection\label{para:VS.B01.P01.C08.S10.SS04.P001}

Kant distinguishes sensibility from understanding.
Through sensibility, objects are given to us in
intuition; through understanding, they are thought
under concepts. Neither faculty can perform the
complete task of cognition in isolation. Sensibility
without conceptual determination does not yield
cognition of an object, while concepts without an
intuition to which they can be applied do not provide
knowledge of an object.
\cite{KantCritiquePureReason}

% Structural ID: VS.B01.P01.C08.S10.SS04.P002
\phantomsection\label{para:VS.B01.P01.C08.S10.SS04.P002}

Space and time are the forms of human sensible intuition.
They are not introduced as empirical properties
abstracted from independently observed spatial and
temporal relations. They are the a priori forms under
which objects can be given to human sensibility.
Consequently, the objects of possible experience
are presented under conditions that belong to the
structure of sensibility.

% Structural ID: VS.B01.P01.C08.S10.SS04.P003
\phantomsection\label{para:VS.B01.P01.C08.S10.SS04.P003}

The categories are pure concepts of the understanding.
Kant organizes them under the headings of quantity,
quality, relation, and modality. They include unity,
plurality, totality, reality, negation, limitation,
substance, causality, community, possibility,
existence, and necessity. Their function is not
to provide a catalogue of independently discovered
properties. They are rules through which the manifold
of intuition can be thought as belonging to objects
of experience.
\cite{KantCritiquePureReason}

% Structural ID: VS.B01.P01.C08.S10.SS04.P004
\phantomsection\label{para:VS.B01.P01.C08.S10.SS04.P004}

The categories acquire objective validity through
their relation to possible experience. Their use is
legitimate where the conditions of sensible intuition
and the synthesis required for experience are satisfied.
An attempt to apply them beyond these conditions
does not automatically produce knowledge of a
supersensible object. The critical limitation is
therefore internal to the justification of their
objective use.

% Structural ID: VS.B01.P01.C08.S10.SS04.P005
\phantomsection\label{para:VS.B01.P01.C08.S10.SS04.P005}

The category of unity is particularly relevant to
the present comparison, but it must be interpreted
within Kant's system. Its role in the synthesis of
a manifold does not establish that Unity, as an
ontological determination, is identical with the
unity of an object in experience. Likewise, the
category of substance cannot be transferred without
qualification into a general ontology of identity.
The transcendental deduction is required to explain
the legitimate application of these categories.


% =============================================================================
% 5. THE TRANSCENDENTAL UNITY OF APPERCEPTION
% =============================================================================

\subsection{The Transcendental Unity of Apperception}
\label{subsec:ch08-kant-apperception}

% Structural ID: VS.B01.P01.C08.S10.SS05.P001
\phantomsection\label{para:VS.B01.P01.C08.S10.SS05.P001}

The principle of transcendental apperception is
expressed in Kant's claim that the "I think" must
be able to accompany all my representations. The
claim does not mean that every representation must
be explicitly noticed by empirical introspection.
It means that representations must be capable of
belonging to one self-consciousness if they are
to be combined as representations in one possible
experience. The unity is a condition of their
being mine in the relevant transcendental sense.
\cite[B131--B132]{KantCritiquePureReason}

% Structural ID: VS.B01.P01.C08.S10.SS05.P002
\phantomsection\label{para:VS.B01.P01.C08.S10.SS05.P002}

Kant distinguishes the analytical unity of
apperception from its synthetic unity. The analytical
unity expresses the identity of the "I think" across
representations. The synthetic unity concerns the
combination of a manifold under one consciousness.
The latter is not an optional addition to an already
completed unity. It is the condition under which
the analytical representation of one and the same
"I" can apply to a manifold of representations.
\cite[B133--B135]{KantCritiquePureReason}

% Structural ID: VS.B01.P01.C08.S10.SS05.P003
\phantomsection\label{para:VS.B01.P01.C08.S10.SS05.P003}

The unity of apperception is therefore not a
substantial object presented through inner sense.
It is the formal unity under which the manifold
can be combined in one consciousness. Kant's
argument does not begin with an observed soul
that persists as an object through time. It begins
with the conditions under which representations
can be united in a single experience. The identity
involved here is consequently not yet a demonstrated
identity of a persisting mental substance.

% Structural ID: VS.B01.P01.C08.S10.SS05.P004
\phantomsection\label{para:VS.B01.P01.C08.S10.SS05.P004}

The distinction between transcendental and empirical
apperception is decisive. Empirical consciousness
is temporally conditioned and concerns particular
states of inner sense. Transcendental apperception
is the necessary unity presupposed by the possibility
of combining representations in one experience.
The former is given as a succession of states;
the latter is a condition of their possible synthesis.
The two must not be treated as different names
for the same empirical phenomenon.

% Structural ID: VS.B01.P01.C08.S10.SS05.P005
\phantomsection\label{para:VS.B01.P01.C08.S10.SS05.P005}

The transcendental unity of apperception has
objective significance because it is related to
the possibility of representing objects. It is
not merely a private feeling of unity. Yet its
objective significance remains tied to the
conditions of possible experience. It does not
by itself provide a description of the intrinsic
constitution of every being independently of
the conditions under which that being could
be experienced by a human subject.


% =============================================================================
% 6. SYNTHESIS, OBJECTIVITY, AND THE UNITY OF EXPERIENCE
% =============================================================================

\subsection{Synthesis, Objectivity, and the Unity of Experience}
\label{subsec:ch08-kant-synthesis}

% Structural ID: VS.B01.P01.C08.S10.SS06.P001
\phantomsection\label{para:VS.B01.P01.C08.S10.SS06.P001}

The manifold of intuition does not become an
object of experience merely by being given.
It must be synthesized according to rules that
permit its elements to be combined in a unified
representation. The understanding supplies the
functions of synthesis through which the manifold
is brought under concepts. Objectivity therefore
involves more than the occurrence of representations:
it requires a rule-governed unity through which
the manifold can be represented as belonging
to an object.
\cite[B137--B143]{KantCritiquePureReason}

% Structural ID: VS.B01.P01.C08.S10.SS06.P002
\phantomsection\label{para:VS.B01.P01.C08.S10.SS06.P002}

The transcendental deduction seeks to justify
the categories by showing that they are necessary
for the possibility of experience. The categories
are not inferred from the repeated observation
of objects. Rather, Kant argues that the unity
of apperception requires a synthesis governed
by concepts that make objective experience possible.
The deduction thus connects the unity of
self-consciousness with the unity of objects
as they can be experienced.

% Structural ID: VS.B01.P01.C08.S10.SS06.P003
\phantomsection\label{para:VS.B01.P01.C08.S10.SS06.P003}

This connection must be interpreted carefully.
The unity of the subject's synthesis and the unity
of the object as represented are related within
the conditions of cognition. It does not follow
that the object is numerically identical with
the act of synthesis, or that its existence is
produced by the subject. Kant's argument concerns
the conditions under which an object can be
represented as an object of experience.

% Structural ID: VS.B01.P01.C08.S10.SS06.P004
\phantomsection\label{para:VS.B01.P01.C08.S10.SS06.P004}

For the Ontology of Viability, the distinction
raises a fundamental question: can the conditions
of objective identification be treated as
conditions of ontological identity? The answer
cannot be assumed in either direction. An
ontology must explain what it means for a being
to possess identity; an epistemology must explain
how that identity can be represented, judged,
or known. Kant's transcendental argument
illuminates the second task and constrains
metaphysical claims, but it does not settle
the first merely by establishing the necessity
of synthesis.

% Structural ID: VS.B01.P01.C08.S10.SS06.P005
\phantomsection\label{para:VS.B01.P01.C08.S10.SS06.P005}

The strongest contribution of Kant's account
is its demonstration that objectivity cannot
be reduced to an unstructured collection of
sensory contents. The representation of an
object presupposes relations of unity and
rule-governed synthesis. The strongest limitation
for the present comparison is that this result
does not yet identify the ontological ground
of the being represented. The explanatory
conditions of objectivity and the ontological
identity of an object remain distinct questions.


% =============================================================================
% 7. APPERCEPTIVE IDENTITY AND THE PROBLEM OF THE SELF
% =============================================================================

\subsection{Apperceptive Identity and the Problem of the Self}
\label{subsec:ch08-kant-self}

% Structural ID: VS.B01.P01.C08.S10.SS07.P001
\phantomsection\label{para:VS.B01.P01.C08.S10.SS07.P001}

Kant's account of apperception contains a claim of
identity: the "I think" must be capable of accompanying
the manifold of representations that can belong
to one consciousness. This identity is not derived
from observing a persisting inner object. It is
presupposed in the synthesis through which the
manifold can be combined. The claim concerns
the formal unity of self-consciousness and the
possibility of representing diverse contents
as belonging to one subject of experience.
\cite[B131--B139]{KantCritiquePureReason}

% Structural ID: VS.B01.P01.C08.S10.SS07.P002
\phantomsection\label{para:VS.B01.P01.C08.S10.SS07.P002}

The identity of apperception must therefore be
distinguished from the identity of an empirical
person considered as an object in time. The
empirical person is encountered through bodily
and psychological states, relations, and changes.
The transcendental unity is not itself given
as one more item in this succession. It is the
formal condition under which the succession
can be combined in one consciousness. The
two uses of identity belong to different
explanatory levels.

% Structural ID: VS.B01.P01.C08.S10.SS07.P003
\phantomsection\label{para:VS.B01.P01.C08.S10.SS07.P003}

This distinction prevents a common inference:
because the "I think" is represented as identical
across a manifold, the subject must therefore
be a simple and enduring substance. Kant denies
that the formal identity of apperception provides
such knowledge. To infer substantiality from
the logical or transcendental unity of the
thinking subject is to move from a condition
of thought to a determination of an object
without supplying the required intuition.

\footnote{For the interpretive difficulties surrounding
Kant's account of the self and the Third Paralogism,
see Karl Ameriks, \textit{Kant's Theory of Mind:
An Analysis of the Paralogisms of Pure Reason}
(Oxford: Clarendon Press, 2000), 128--176.}

% Structural ID: VS.B01.P01.C08.S10.SS07.P004
\phantomsection\label{para:VS.B01.P01.C08.S10.SS07.P004}

The argument is particularly important because
it blocks both an empiricist reduction and a
rationalist overreach. The unity of the subject
cannot be explained simply as an observed
collection of changing mental states; but neither
can the formal unity of thought be converted
into knowledge of a metaphysical soul. Kant
therefore preserves a distinction between the
unity required for experience and the substantial
constitution of the subject as an object.

% Structural ID: VS.B01.P01.C08.S10.SS07.P005
\phantomsection\label{para:VS.B01.P01.C08.S10.SS07.P005}

The distinction also bears on the relation
between identity and continuity. The fact that
a manifold can be unified under one apperception
does not establish that every empirical state
is continuously conscious, nor does it provide
a complete criterion for personal persistence
through every interruption or transformation.
The transcendental condition of unity and the
temporal history of a person must be examined
separately.

% Structural ID: VS.B01.P01.C08.S10.SS07.P006
\phantomsection\label{para:VS.B01.P01.C08.S10.SS07.P006}

The following diagram summarizes the role of
synthesis and apperceptive unity in Kant's
account. Its purpose is to distinguish the
conditions under which a manifold can be
represented as one experience from the
metaphysical conclusions that cannot be
inferred from that unity alone.

% =============================================================================
% Project : THE PHILOSOPHY OF VIABILITY
% Book    : Book I — Ontology of Viability
% Book ID : VS.B01
% Part    : Part I — Identity and Manifestation
% Part ID : VS.B01.P01
% Chapter : Chapter 08 — Comparative Philosophical Examination
% Chapter ID : VS.B01.P01.C08
%
% File :
%   4_ASSETS/diagrams/
%   VS.B01.P01.C08.S10.SS07.P006.DG01.kant_synthesis_and_apperceptive_unity.tex
%
% Section :
%   Section 10 — Immanuel Kant
%
% Asset :
%   Diagram DG01 — Kant's Synthesis and Apperceptive Unity
%
% Structural Identifier :
%   VS.B01.P01.C08.S10.SS07.P006.DG01
%
% Asset Type : DG
%
% Purpose :
%   To distinguish the manifold of intuition, synthesis,
%   the transcendental unity of apperception, and the
%   representation of an object in possible experience.
%
% Intellectual Benefit :
%   The diagram makes visible the structure of Kant's
%   transcendental account and clarifies the relation
%   between apperceptive unity and objectivity.
%
% Aesthetic / Experiential Benefit :
%   It provides a visual pause within a conceptually dense
%   discussion and gives the reader a compact overview of
%   the central argument.
%
% Representational Limit :
%   The diagram represents the conditions of possible
%   experience in Kant's account. It does not depict an
%   ontological production of beings by the subject.
%   The lower distinction explicitly marks the limit
%   of the inference from apperceptive unity to a
%   substantial or numerically persistent soul.
%
% Dependencies :
%   TikZ; positioning library.
%   Both are used elsewhere in the project.
%
% =============================================================================

\begin{figure}[htbp]
    \centering

    \begin{tikzpicture}[
        font=\small,
        node distance=0.68cm,
        process/.style={
            draw=blue!55!black,
            fill=blue!6,
            rounded corners=3pt,
            line width=0.7pt,
            text width=9.0cm,
            minimum height=0.82cm,
            align=center,
            inner sep=6pt
        },
        central/.style={
            process,
            fill=blue!15,
            font=\small\bfseries
        },
        result/.style={
            draw=teal!55!black,
            fill=teal!7,
            rounded corners=3pt,
            line width=0.7pt,
            text width=9.0cm,
            minimum height=0.82cm,
            align=center,
            inner sep=6pt
        },
        limit/.style={
            draw=black!55,
            fill=black!4,
            rounded corners=3pt,
            line width=0.7pt,
            text width=11.0cm,
            align=center,
            font=\footnotesize,
            inner sep=7pt
        },
        flow/.style={
            ->,
            draw=black!65,
            line width=0.8pt
        },
        distinction/.style={
            draw=red!55!black,
            fill=red!4,
            rounded corners=3pt,
            line width=0.7pt,
            text width=11.0cm,
            align=center,
            font=\footnotesize,
            inner sep=7pt
        }
    ]

    \node[process] (manifold) {
        \textbf{Manifold of Sensible Intuition}\\
        Representations given under the forms of space and time
    };

    \node[process, below=of manifold] (synthesis) {
        \textbf{Synthesis of the Manifold}\\
        Combination according to rules of the understanding
    };

    \node[central, below=of synthesis] (apperception) {
        \textbf{Transcendental Unity of Apperception}\\
        The manifold must be capable of belonging to one
        self-consciousness: the unity of the ``I think''
    };

    \node[result, below=of apperception] (categories) {
        \textbf{Categories and Objective Unity}\\
        The manifold is represented under rules that make
        the cognition of objects possible
    };

    \node[result, below=of categories] (experience) {
        \textbf{Unified Possible Experience}\\
        Objects can be represented within one rule-governed
        experience
    };

    \draw[flow] (manifold) -- (synthesis);
    \draw[flow] (synthesis) -- (apperception);
    \draw[flow] (apperception) -- (categories);
    \draw[flow] (categories) -- (experience);

    \node[limit, below=0.85cm of experience] (limit) {
        \textbf{Critical limit of the inference}\\
        The unity required for objective experience does not,
        by itself, establish that the transcendental subject
        is a simple, enduring, substantial soul.
    };

    \node[distinction, below=0.55cm of limit] (ontology) {
        \textbf{Ontological question left open}\\
        What constitutes the identity of a being independently
        of the conditions under which it is represented as
        an object of possible experience?
    };

    \end{tikzpicture}

    \caption{Kant's account of synthesis, transcendental
    apperception, and the unity of possible experience.
    The diagram distinguishes the conditions of objective
    cognition from the ontological constitution of the
    subject or of the objects represented.}

    \label{fig:VS.B01.P01.C08.S10.SS07.P006.DG01}
\end{figure}

% =============================================================================
% END OF DIAGRAM DG01
% =============================================================================


% =============================================================================
% 8. THE PARALOGISMS AND THE LIMITS OF RATIONAL PSYCHOLOGY
% =============================================================================

\subsection{The Paralogisms and the Limits of Rational Psychology}
\label{subsec:ch08-kant-paralogisms}

% Structural ID: VS.B01.P01.C08.S10.SS08.P001
\phantomsection\label{para:VS.B01.P01.C08.S10.SS08.P001}

The Paralogisms examine the attempts of rational
psychology to derive the nature of the soul from
the bare representation "I think." Kant argues
that the subject of thought is treated as though
it were an object about which the categories
could yield determinate knowledge, even though
the required intuition of such an object has
not been supplied. The error is not merely
a faulty empirical observation. It concerns
the illegitimate application of concepts beyond
the conditions that authorize their objective use.
\cite[A341/B399--A405/B432]{KantCritiquePureReason}

% Structural ID: VS.B01.P01.C08.S10.SS08.P002
\phantomsection\label{para:VS.B01.P01.C08.S10.SS08.P002}

The rational psychologist attempts to infer
that the soul is a substance, simple, numerically
identical, and distinct from the body. Kant
examines the inferential transitions involved
in these claims. The representation of the
subject as the logical subject of judgments
does not prove that it is a substance in the
ontological sense. The unity of thought does
not establish simplicity as a property of
an independently given object. The identity
of the "I" in thought does not prove the
persistence of a person as a substantial
being through time.

% Structural ID: VS.B01.P01.C08.S10.SS08.P003
\phantomsection\label{para:VS.B01.P01.C08.S10.SS08.P003}

Kant's criticism is directed against a
misapplication of the categories. The category
of substance has legitimate use in the
determination of appearances within possible
experience. It cannot be applied to the
transcendental subject merely because the
subject functions as the logical subject
of thought. The same applies to the inference
from the simplicity of the representation
"I" to the simplicity of a soul as an object.

% Structural ID: VS.B01.P01.C08.S10.SS08.P004
\phantomsection\label{para:VS.B01.P01.C08.S10.SS08.P004}

The critical result is a limitation on
theoretical knowledge, not a demonstration
that no subject exists or that the person
is nothing but a succession of representations.
Kant rejects the claim that pure rational
psychology can establish the substantial
constitution of the soul from the formal
representation of self-consciousness. The
failure of that proof must not be confused
with a proof of the opposite ontology.

% Structural ID: VS.B01.P01.C08.S10.SS08.P005
\phantomsection\label{para:VS.B01.P01.C08.S10.SS08.P005}

This point is especially relevant to the
Ontology of Viability. The failure of a
particular inference from self-consciousness
to substance does not establish that identity
is reducible to psychological synthesis.
It leaves open the ontological question
concerning the identity of the subject and
the conditions under which it persists.
Kant's criticism therefore blocks an
unwarranted metaphysical inference while
leaving room for an independent ontological
investigation.


% =============================================================================
% 9. SUBSTANCE, PERMANENCE, AND THE FIRST ANALOGY
% =============================================================================

\subsection{Substance, Permanence, and the First Analogy}
\label{subsec:ch08-kant-substance}

% Structural ID: VS.B01.P01.C08.S10.SS09.P001
\phantomsection\label{para:VS.B01.P01.C08.S10.SS09.P001}

Kant's First Analogy of Experience states that
substance is permanent in all change of
appearances. The argument concerns the
conditions under which temporal relations
such as succession and simultaneity can be
determined objectively. Since time itself
is not perceived as an object, Kant argues
that the determination of changes requires
a persistent correlate in appearances.
The principle is therefore presented as
a condition of temporal experience.
\cite[A182/B224--A189/B232]{KantCritiquePureReason}

% Structural ID: VS.B01.P01.C08.S10.SS09.P002
\phantomsection\label{para:VS.B01.P01.C08.S10.SS09.P002}

The permanence of substance is not simply
a generalization from the observation that
some objects happen to remain unchanged.
Kant intends it as a synthetic a priori
principle governing the temporal determination
of appearances. Change is represented as
a determination of what persists; without
a persistent reference, the temporal relation
between changing states could not be
objectively established in the manner
required by experience.

% Structural ID: VS.B01.P01.C08.S10.SS09.P003
\phantomsection\label{para:VS.B01.P01.C08.S10.SS09.P003}

The principle nevertheless has a restricted
domain. It concerns substance as appearance
within possible experience. It does not
authorize a direct inference from phenomenal
permanence to the constitution of things
considered independently of the conditions
of experience. Nor does it establish that
every particular object or person remains
numerically identical merely because a
persistent substrate is required for the
objective determination of change.

% Structural ID: VS.B01.P01.C08.S10.SS09.P004
\phantomsection\label{para:VS.B01.P01.C08.S10.SS09.P004}

A further distinction is required between
permanence and identity. Kant's principle
establishes the role of permanence in the
temporal determination of appearances.
It does not follow that permanence is
identical with identity, or that identity
can be defined as the persistence of an
unchanging substrate. The principle concerns
a specific condition of temporal experience;
the ontological question of what makes
a being the being that it is requires
a distinct analysis.

% Structural ID: VS.B01.P01.C08.S10.SS09.P005
\phantomsection\label{para:VS.B01.P01.C08.S10.SS09.P005}

The comparison with the preceding discussion
of Hume is instructive. Hume examines the
origin of the idea of substance and the
psychological attribution of identity to
changing or related perceptions. Kant
asks what must be presupposed for objective
temporal experience and argues for the
necessity of permanence in appearances.
The two projects therefore address different
questions, even where they employ related
terms. Kant's transcendental argument is
not merely a reversal of Hume's psychology;
it concerns the objective conditions of
experience.
\cite{KantCritiquePureReason}

% Structural ID: VS.B01.P01.C08.S10.SS09.P006
\phantomsection\label{para:VS.B01.P01.C08.S10.SS09.P006}

The Ontology of Viability can accept the
importance of distinguishing persistence,
continuity, and identity without adopting
Kant's principle as a universal ontological
definition. The question whether a particular
being persists, whether its manifestations
continue, and whether its identity remains
the same must be addressed according to
the relevant ontological conditions. A
transcendental requirement for the temporal
determination of appearances cannot by
itself settle every such case.


% =============================================================================
% 10. PERSONAL IDENTITY AND THE THIRD PARALOGISM
% =============================================================================

\subsection{Personal Identity and the Third Paralogism}
\label{subsec:ch08-kant-personal-identity}

% Structural ID: VS.B01.P01.C08.S10.SS10.P001
\phantomsection\label{para:VS.B01.P01.C08.S10.SS10.P001}

Kant's treatment of personal identity in the
Paralogisms is directed against the inference
that the identity of the "I" in thought
establishes the numerical identity of a
person as a persisting substance. The
representation that I am conscious of
myself as the same subject in different
acts of thought does not, without further
argument, establish that the same substantial
being persists through all changes of
inner states. The identity of apperception
and the identity of a person as an object
are not equivalent claims.
\cite[A361/B408--A363/B411]{KantCritiquePureReason}

% Structural ID: VS.B01.P01.C08.S10.SS10.P002
\phantomsection\label{para:VS.B01.P01.C08.S10.SS10.P002}

Kant's criticism is particularly important
because it separates the formal identity
of the thinking subject from the empirical
and metaphysical claims associated with
personal persistence. The subject may be
represented as identical in the unity of
consciousness, but the intuition required
to determine the subject as an enduring
object is not supplied by the mere
representation "I think." The inference
to a persisting soul therefore exceeds
what the formal unity alone can establish.

% Structural ID: VS.B01.P01.C08.S10.SS10.P003
\phantomsection\label{para:VS.B01.P01.C08.S10.SS10.P003}

This analysis must not be interpreted as
a straightforward denial of personal
identity. Kant's argument concerns the
limits of theoretical proof. It does not
demonstrate that a person lacks identity,
nor does it replace the person with a
bundle of perceptions. The conclusion
is that the identity of the subject as
represented in apperception cannot be
used as sufficient evidence for the
substantial persistence of the soul.

% Structural ID: VS.B01.P01.C08.S10.SS10.P004
\phantomsection\label{para:VS.B01.P01.C08.S10.SS10.P004}

The difference from Locke and Hume is
significant. Locke connects personal
identity with continuity of consciousness
and the appropriation of past actions.
Hume examines the succession of perceptions
and the operations through which identity
is attributed. Kant asks whether the
formal identity presupposed in self-conscious
thought can establish the identity of a
person as an enduring substance. These
are related but distinct explanatory
strategies. The present comparison must
not treat them as competing answers to
one undifferentiated question.
\cite{KantCritiquePureReason}

% Structural ID: VS.B01.P01.C08.S10.SS10.P005
\phantomsection\label{para:VS.B01.P01.C08.S10.SS10.P005}

The unresolved ontological issue is what
grounds the identity of a person when
the formal unity of apperception is
distinguished from empirical continuity
and from the substantiality of the soul.
Kant establishes a limit on one route
to knowledge of personal identity. He
does not thereby provide a complete
positive ontology of the person. This
limitation is not a defect in the
specific argument against rational
psychology; it marks the boundary
between the task Kant undertakes and
the further ontological task pursued
in this book.


% =============================================================================
% 11. PHENOMENA, NOUMENA, AND THE LIMITS OF ONTOLOGICAL INFERENCE
% =============================================================================

\subsection{Phenomena, Noumena, and the Limits of Ontological Inference}
\label{subsec:ch08-kant-phenomena-noumena}

% Structural ID: VS.B01.P01.C08.S10.SS11.P001
\phantomsection\label{para:VS.B01.P01.C08.S10.SS11.P001}

Kant's distinction between appearances and things
considered in themselves marks a boundary within
the conditions of human cognition. Appearances
are objects insofar as they are given under
the forms of human sensibility and thought
under the categories. Things considered in
themselves are considered independently of
those conditions. The distinction does not
permit an unrestricted description of a
second world populated by objects whose
properties can be specified through the
same categories used in experience.
\cite[A235/B294--A260/B315]{KantCritiquePureReason}

% Structural ID: VS.B01.P01.C08.S10.SS11.P002
\phantomsection\label{para:VS.B01.P01.C08.S10.SS11.P002}

The term noumenon has a limiting function
in Kant's critical philosophy. It indicates
the concept of an object considered apart
from sensible intuition, but such a concept
does not provide theoretical knowledge of
a supersensible object. The distinction
therefore restricts the scope of claims
that can be made through the categories.
The absence of sensible conditions does
not become a licence to attribute determinate
properties to an unknowable object.

% Structural ID: VS.B01.P01.C08.S10.SS11.P003
\phantomsection\label{para:VS.B01.P01.C08.S10.SS11.P003}

The interpretation of the thing in itself
remains contested. Some readings emphasize
its realist implications, while others
stress its function as a limiting concept
within transcendental idealism. The present
section does not resolve this interpretive
dispute by assuming that Kant posits a
fully describable noumenal realm. It
distinguishes the critical restriction
on knowledge from stronger metaphysical
claims that require independent textual
and philosophical argument.

\footnote{The relation between transcendental idealism,
the status of things in themselves, and the limits
of theoretical cognition is a major interpretive
problem in Kant scholarship. See Allison,
\textit{Kant's Transcendental Idealism}, 2nd ed.}

% Structural ID: VS.B01.P01.C08.S10.SS11.P004
\phantomsection\label{para:VS.B01.P01.C08.S10.SS11.P004}

For the Ontology of Viability, the distinction
raises a serious challenge: if the identity
of a being is not simply identical with
its experienced manifestations, what
warrants claims about that identity beyond
what is manifested? The answer cannot
be to declare identity unknowable merely
because it is not reducible to an appearance.
Nor can the answer be to treat identity
as a noumenal object whose properties
are known independently of any argument.
The ontological status of identity must
be established through the foundations
and arguments of Book I, not by importing
Kant's distinction as a ready-made solution.

% Structural ID: VS.B01.P01.C08.S10.SS11.P005
\phantomsection\label{para:VS.B01.P01.C08.S10.SS11.P005}

The comparison therefore remains asymmetrical.
Kant's critical distinction constrains
the extension of theoretical knowledge
beyond possible experience. The Ontology
of Viability investigates the relation
between Being, Identity, and Manifestation
as ontological questions. These projects
may intersect in their concern with the
limits of determination, but their
explanatory aims and starting points
must not be conflated.


% =============================================================================
% 12. IDENTITY, DIFFERENCE, AND TRANSCENDENTAL REFLECTION
% =============================================================================

\subsection{Identity, Difference, and Transcendental Reflection}
\label{subsec:ch08-kant-identity-difference}

% Structural ID: VS.B01.P01.C08.S10.SS12.P001
\phantomsection\label{para:VS.B01.P01.C08.S10.SS12.P001}

Kant's discussion of the Amphiboly of the Concepts
of Reflection examines the conditions under which
representations are compared as identical or
different, in agreement or opposition. He argues
that logical comparison alone is insufficient
when the relevant representations may belong
to different cognitive faculties or occupy
different relations to sensibility and understanding.
Transcendental reflection is required to determine
the standpoint from which the comparison is made.
\cite[A260/B316--A292/B349]{KantCritiquePureReason}

% Structural ID: VS.B01.P01.C08.S10.SS12.P002
\phantomsection\label{para:VS.B01.P01.C08.S10.SS12.P002}

The example of two drops of water illustrates
the difference between conceptual indiscernibility
and numerical identity. Two objects may be
qualitatively indistinguishable while remaining
numerically distinct because they occupy
different positions in space. The comparison
shows that the identity of a concept and the
identity of an object are not interchangeable.
The conditions under which objects are given
must be considered when numerical identity
is judged.

% Structural ID: VS.B01.P01.C08.S10.SS12.P003
\phantomsection\label{para:VS.B01.P01.C08.S10.SS12.P003}

This analysis is relevant to the relation
between identity and distinction. Kant
demonstrates that a purely logical comparison
of concepts does not settle every question
about the identity or difference of objects.
The conditions of sensible presentation
matter to the objective determination of
numerical difference. Yet this argument
does not provide a complete ontology of
identity; it clarifies the conditions
under which identity and difference
can be judged within experience.

% Structural ID: VS.B01.P01.C08.S10.SS12.P004
\phantomsection\label{para:VS.B01.P01.C08.S10.SS12.P004}

The Ontology of Viability can recognize
the importance of distinguishing identity,
qualitative indiscernibility, numerical
identity, and difference without adopting
Kant's transcendental framework as its
own foundation. The question whether two
manifestations belong to one identity,
or whether two qualitatively similar
beings are numerically distinct, requires
an ontological analysis of the relevant
being and its determinations. The
conditions of comparison are important,
but they do not substitute for that analysis.

% Structural ID: VS.B01.P01.C08.S10.SS12.P005
\phantomsection\label{para:VS.B01.P01.C08.S10.SS12.P005}

The distinction also cautions against
treating identity as a simple repetition
of predicates. A comparison of descriptions
may establish similarity or indiscernibility
under a specified respect, but it does
not automatically settle whether the
entities compared are one and the same.
Kant's analysis therefore supports
conceptual discipline in the treatment
of identity, while leaving the ontological
ground of identity open to further inquiry.


% =============================================================================
% 13. THE STRONGEST KANTIAN CHALLENGE TO THE ONTOLOGY OF VIABILITY
% =============================================================================

\subsection{The Strongest Kantian Challenge to the Ontology of Viability}
\label{subsec:ch08-kant-challenge}

% Structural ID: VS.B01.P01.C08.S10.SS13.P001
\phantomsection\label{para:VS.B01.P01.C08.S10.SS13.P001}

A strong Kantian challenge begins by asking how
the Ontology of Viability can claim to describe
identity independently of the conditions under
which identity is determined, represented, or
known. If every philosophical assertion about
an object is made through concepts and judgments,
then the ontology must explain how its claims
avoid treating the conditions of possible
experience as though they were transparent
descriptions of things independently of those
conditions. The challenge is not that ontology
is impossible, but that its claims require
an account of their scope and warrant.

% Structural ID: VS.B01.P01.C08.S10.SS13.P002
\phantomsection\label{para:VS.B01.P01.C08.S10.SS13.P002}

The challenge can be sharpened in relation
to the distinction between Identity and
Manifestation. If identity is not itself
one more manifestation, how can an ontology
describe its relation to manifestations
without converting it into an inaccessible
substrate? Conversely, if identity is
known only through manifestations, how
can it retain an ontological status that
is not reducible to the manifested
properties or to the conditions of
their apprehension? Kant's distinction
between appearance and thing in itself
makes this tension especially visible.

% Structural ID: VS.B01.P01.C08.S10.SS13.P003
\phantomsection\label{para:VS.B01.P01.C08.S10.SS13.P003}

A second challenge concerns the unity of
the object. Kant may ask whether the
Ontology of Viability has supplied a
sufficient account of the conditions
under which a plurality of determinations
and manifestations belongs to one being.
The mere assertion that identity is
ontologically inseparable from Being
does not yet explain how the unity of
a being is distinguished from the unity
through which an observer represents
a manifold as one object. The ontology
must not rely on a cognitive act to
establish the unity it claims to describe.

% Structural ID: VS.B01.P01.C08.S10.SS13.P004
\phantomsection\label{para:VS.B01.P01.C08.S10.SS13.P004}

A third challenge concerns the status of
the categories and distinctions employed
by the ontology. If its concepts are
used to make claims about Being, Identity,
Determination, and Manifestation, it
must explain whether these concepts
describe ontological relations, organize
philosophical analysis, or perform both
functions in carefully distinguished
ways. Without such an account, the
ontology risks confusing its conceptual
architecture with the structure of
what it investigates.

% Structural ID: VS.B01.P01.C08.S10.SS13.P005
\phantomsection\label{para:VS.B01.P01.C08.S10.SS13.P005}

The strongest form of the Kantian challenge
therefore demands a clear separation
between the conditions of being, the
conditions of manifestation, and the
conditions of knowing or judging what
is manifested. It asks the Ontology of
Viability to show that its ontological
commitments are not covertly derived
from the structure of cognition. This
is a substantive challenge and cannot
be answered merely by insisting that
Book I is an ontology rather than an
epistemology.


% =============================================================================
% 14. BILATERAL EVALUATION AND CLAIM STATUS
% =============================================================================

\subsection{Bilateral Evaluation and Claim Status}
\label{subsec:ch08-kant-evaluation}

% Structural ID: VS.B01.P01.C08.S10.SS14.P001
\phantomsection\label{para:VS.B01.P01.C08.S10.SS14.P001}

Kant's strongest contribution to the present
comparison is his demonstration that the
unity of objective experience requires
conditions that cannot be derived from
the mere accumulation of unconnected
representations. The transcendental unity
of apperception and the synthesis of the
manifold provide an account of how
representations can belong to one
experience and how objects can be
represented under necessary rules.
This is a significant analysis of the
conditions of objectivity.

% Structural ID: VS.B01.P01.C08.S10.SS14.P002
\phantomsection\label{para:VS.B01.P01.C08.S10.SS14.P002}

The comparison also requires a distinction between
the explanatory success of a transcendental argument
and the ontological scope of its conclusion. An argument
may establish that a condition is necessary for the
possibility of experience without establishing that
the condition exhausts the constitution of the beings
experienced. This distinction governs the evaluation
of each of the claims summarized in the following table.

% Structural ID: VS.B01.P01.C08.S10.SS14.P003
\phantomsection\label{para:VS.B01.P01.C08.S10.SS14.P003}

The following table distinguishes Kant's
principal claims, the explanatory work
they perform, and the conclusions that
cannot be inferred from them without
additional argument.

% =============================================================================
% Project : THE PHILOSOPHY OF VIABILITY
% Book    : Book I — Ontology of Viability
% Book ID : VS.B01
% Part    : Part I — Identity and Manifestation
% Part ID : VS.B01.P01
% Chapter : Chapter 08 — Comparative Philosophical Examination
% Chapter ID : VS.B01.P01.C08
%
% File :
%   4_ASSETS/tables/
%   VS.B01.P01.C08.S10.SS14.P003.TL01.kant_claims_and_ontological_limits.tex
%
% Section :
%   Section 10 — Immanuel Kant
%
% Asset :
%   Table TL01 — Kant's Claims and Their Ontological Limits
%
% Structural Identifier :
%   VS.B01.P01.C08.S10.SS14.P003.TL01
%
% Asset Type : TL
%
% Purpose :
%   To distinguish the principal claims in Kant's critical
%   philosophy according to their function, explanatory
%   reach, and limits for an ontological examination of
%   identity.
%
% Intellectual Benefit :
%   The table prevents the conflation of apperceptive unity,
%   empirical self-consciousness, substantiality, permanence,
%   and personal identity.
%
% Aesthetic / Experiential Benefit :
%   It provides a compact comparative overview after the
%   extended argument and reduces the cognitive burden of
%   retaining several closely related Kantian distinctions.
%
% Representational Limit :
%   The table is an analytical summary. It does not imply
%   that Kant's critical philosophy is internally uncontested
%   or that the comparative limits listed here settle every
%   interpretive dispute.
%
% =============================================================================

\begin{center}
\small
\renewcommand{\arraystretch}{1.18}
\setlength{\tabcolsep}{4pt}

\begin{longtable}{
    >{\raggedright\arraybackslash}p{2.3cm}
    >{\raggedright\arraybackslash}p{0.27\textwidth}
    >{\raggedright\arraybackslash}p{0.27\textwidth}
    >{\raggedright\arraybackslash}p{0.27\textwidth}
}

\caption{Kant's Principal Claims and Their Ontological Limits}
\label{tab:VS.B01.P01.C08.S10.SS14.P003.TL01} \\

\toprule
\textbf{Problem}
&
\textbf{Kant's account}
&
\textbf{What it establishes}
&
\textbf{What it does not establish}
\\
\midrule
\endfirsthead

\toprule
\textbf{Problem}
&
\textbf{Kant's account}
&
\textbf{What it establishes}
&
\textbf{What it does not establish}
\\
\midrule
\endhead

\midrule
\multicolumn{4}{r}{\footnotesize\textit{Continued on next page}}\\
\endfoot

\bottomrule
\endlastfoot

Copernican turn
&
Objects, insofar as they are objects of possible
experience, conform to a priori conditions of
human cognition.
&
The transcendental conditions under which
objective experience is possible.
&
That the existence or ontological constitution
of beings is produced by the subject.
\\

\addlinespace

Categories
&
Pure concepts of the understanding provide
rules for synthesizing the manifold of intuition.
&
The conceptual conditions under which objects
can be thought in possible experience.
&
That the categories constitute an exhaustive
inventory of the determinations of Being.
\\

\addlinespace

Transcendental apperception
&
The manifold must be capable of belonging
to one self-consciousness under the unity
of the ``I think.''
&
The formal unity required for the synthesis
of representations in one experience.
&
That the subject is thereby known as a
simple or enduring substance.
\\

\addlinespace

Transcendental deduction
&
The categories are justified through their
necessary relation to the possibility of
experience.
&
The objective validity of the categories
within their legitimate domain.
&
Their unrestricted application beyond
possible experience.
\\

\addlinespace

First Analogy
&
Substance is permanent in all change
of appearances.
&
A principle for the objective temporal
determination of appearances.
&
That every particular being or person
is thereby proved numerically identical
through all transformations.
\\

\addlinespace

Paralogisms
&
The formal unity of the thinking subject
does not prove the substantiality,
simplicity, or persistence of the soul.
&
The limits of rational psychology's
inferences from the ``I think.''
&
That the subject does not exist or
that personal identity is merely
a bundle of perceptions.
\\

\addlinespace

Personal identity
&
The identity of the ``I'' in thought
cannot alone establish the identity
of a person as a persisting substance.
&
A distinction between apperceptive
identity and substantial personal
persistence.
&
A complete positive ontology of
personal identity.
\\

\addlinespace

Phenomena and noumena
&
The categories have objective use
only under the conditions of possible
experience.
&
A critical boundary on theoretical
knowledge.
&
A determinate description of things
in themselves as objects of
theoretical knowledge.
\\

\addlinespace

Identity and difference
&
Objective comparison requires
transcendental reflection on the
relation of representations to
the faculties of cognition.
&
The insufficiency of merely logical
comparison for every judgment
about objects.
&
A complete account of the ontological
ground of identity and distinction.
\\

\end{longtable}

\vspace{0.5em}

\begin{minipage}{0.95\textwidth}
\footnotesize
\textit{Note:}
The table distinguishes the scope of Kant's transcendental
arguments from the broader ontological conclusions that
might be drawn from them. Such conclusions require
additional argument and must not be treated as automatic
consequences of the critical philosophy.
\end{minipage}

\end{center}

% =============================================================================
% END OF TABLE TL01
% =============================================================================

% Structural ID: VS.B01.P01.C08.S10.SS14.P004
\phantomsection\label{para:VS.B01.P01.C08.S10.SS14.P004}

The limitation of Kant's account, from the
standpoint of this book, is not that
transcendental conditions are irrelevant
to ontology. It is that a demonstration
of the conditions under which objects
can be experienced does not by itself
establish the complete ontological
constitution of those objects. The
conditions of objectivity and the
identity of the beings represented
must be related through further
argument rather than identified.

% Structural ID: VS.B01.P01.C08.S10.SS14.P005
\phantomsection\label{para:VS.B01.P01.C08.S10.SS14.P005}

The Ontology of Viability accepts the
need to distinguish ontological identity
from the cognitive conditions under
which identity is recognized. It also
accepts the force of Kant's warning
against deriving the substantiality
of the self from the formal unity
of self-consciousness. These points
support the methodological independence
of the ontological inquiry. They do
not, however, establish that the
ontology must adopt transcendental
idealism or treat the categories as
the ultimate conditions of Being.

% Structural ID: VS.B01.P01.C08.S10.SS14.P006
\phantomsection\label{para:VS.B01.P01.C08.S10.SS14.P006}

The comparison remains open at one
important point. Kant's critical
philosophy restricts claims about
objects beyond possible experience.
The Ontology of Viability must
explain how its claims about Being
and Identity are justified without
turning them into unargued descriptions
of an unknowable substrate. The
present section identifies this
challenge but does not claim to
have resolved every issue concerning
the relation between ontology and
the limits of cognition.

% Structural ID: VS.B01.P01.C08.S10.SS14.P007
\phantomsection\label{para:VS.B01.P01.C08.S10.SS14.P007}

The claim status of the principal
conclusions must remain explicit.
The historical reconstruction of
Kant's distinction between apperceptive
unity and substantial self-knowledge
is textually supported. The claim
that this distinction does not
establish the ontological identity
of beings is an argued comparative
conclusion. The adequacy of the
Ontology of Viability's own account
of identity under Kant's challenge
remains a matter for continuing
argument. It is not treated here
as a completed proof.

% Structural ID: VS.B01.P01.C08.S10.SS14.P008
\phantomsection\label{para:VS.B01.P01.C08.S10.SS14.P008}

The functional form and epistemic status
of a claim must also remain distinct.
A historical thesis about Kant's
position is not thereby a foundational
commitment of Book I. A proposition
about the limits of Kant's inference
is not an axiom of the Ontology of
Viability. A hypothesis concerning
the ontological relation between
identity and manifestation remains
a hypothesis even when it is
suggested by the comparison.


% =============================================================================
% 15. TRANSITION
% =============================================================================

\subsection{Transition}
\label{subsec:ch08-kant-transition}

% Structural ID: VS.B01.P01.C08.S10.SS15.P001
\phantomsection\label{para:VS.B01.P01.C08.S10.SS15.P001}

Kant's critical philosophy establishes
a demanding account of the conditions
under which a manifold can be united
as objective experience. It also
shows why the formal identity of
self-consciousness cannot be converted
without further argument into knowledge
of a simple, enduring, substantial soul.
His analysis of permanence, personal
identity, and the limits of rational
psychology therefore provides an
important test for any ontology that
distinguishes the identity of a being
from the ways in which it is manifested
and cognized.

% Structural ID: VS.B01.P01.C08.S10.SS15.P002
\phantomsection\label{para:VS.B01.P01.C08.S10.SS15.P002}

The comparison has not established
that Kant's transcendental conditions
are the ontological ground of identity.
Nor has it established that the
Ontology of Viability can disregard
the problem of the conditions under
which identity is determined and
judged. The remaining task is to
examine how other philosophical
traditions approach the relation
between Being, Identity, determination,
and manifestation, while preserving
the distinction between ontological
claims and the conditions of their
intelligibility.

% Structural ID: VS.B01.P01.C08.S10.SS15.P003
\phantomsection\label{para:VS.B01.P01.C08.S10.SS15.P003}

The next comparative examination
continues the inquiry into the
ontological status of identity and
the limits of its determination.
The questions raised by Kant will
remain relevant wherever a philosophical
system distinguishes what is from
the conditions under which it can
be experienced, represented, or known.


% =============================================================================
% END OF SECTION 10 OF CHAPTER 08
% =============================================================================